\unit{Thuyết trình 236 -- Các lược đồ Picard: các tính chất tổng quát}
\sid{236}
\markboth{Thuyết trình 236}{Thuyết trình 236}
\src{221}
\seminar{Năm thứ 14\textsuperscript{}, 1961/62, số 236\textsuperscript{}\\Tháng Năm 1962}
\paperhead{KỸ THUẬT HẠ XUỐNG VÀ CÁC ĐỊNH LÝ TỒN TẠI TRONG HÌNH HỌC ĐẠI SỐ\\
VI. CÁC LƯỢC ĐỒ PICARD: CÁC TÍNH CHẤT TỔNG QUÁT}%
  {do Alexander GROTHENDIECK}

\fsection{0}{Các bổ sung cho thuyết trình trước ([3], V)\texorpdfstring{\textsuperscript{(1)}}{(1)}}

Đã có một số tiến bộ liên quan đến các vấn đề tồn tại của những tiền lược đồ
Picard được nêu ra trong [3], V:

\medskip
\noindent a. (Mumford). Nói chung, không đúng rằng khi
$f:X\longrightarrow S$ là một cấu xạ xạ ảnh và khả ly (= phẳng với các thớ
khả ly), tiền lược đồ Picard $\operatorname{Pic}_{X/S}$ tồn tại, ngay cả khi
các thớ của $f$ có chiều 1 và $S$ là phổ của một vành định giá rời rạc đầy đủ.
Một ví dụ thu được bằng cách lấy $S=\operatorname{Spec}\mathbf R[[t]]$ và lấy
$X$ là lược đồ con của $\mathbf P^2_S$ được xác định bởi phương trình (trong đó
$x,y,z$ là các biến thuần nhất)
\[
x^2+y^2=tz^2.
\]
Như vậy, đây là một conic suy biến về mặt hình học thành hai đường thẳng giao
nhau, nhưng thớ đặc biệt trên trường $\mathbf R$ vẫn bất khả quy (nó được cho
bởi phương trình $x^2+y^2=0$ trên $\mathbf R$). Dễ thấy rằng, sau mở rộng étale
$S'\longrightarrow S$, với $S'=\operatorname{Spec}\mathbf C[[t]]$, tiền lược
đồ Picard của $X'/S'$ tồn tại và có một mô tả tường minh như tổng rời nhau của
các bản sao $\widetilde S'$, trong đó $\widetilde S'$ thu được từ $S'$ bằng
cách tách đôi gốc vô hạn lần. Ta dễ dàng nhận thấy dữ liệu hạ xuống trên
$\operatorname{Pic}_{X'/S'}$ đối với $S'\longrightarrow S$ (ở đây được cho
bởi tác động của nhóm Galois $\mathbf Z/2\mathbf Z$ của $S'$ trên $S$) không
hữu hiệu: nhóm này hoán đổi một số điểm đã tách đôi, nên có những quỹ đạo không
nằm trong một tập mở afin. Tuy nhiên, Mumford có thể chứng minh rằng nếu
$f:X\longrightarrow S$ là một cấu xạ xạ ảnh và khả ly sao cho, với mọi
$s\in S$, các thành phần bất khả quy của $X_s$ là bất khả quy hình học trên
$k(s)$, thì $\operatorname{Pic}_{X/S}$ tồn tại; chứng minh của ông dựa trên
một sự tinh chỉnh định lý thương của mình, xem [7]. Mặt khác, cần lưu ý rằng,
ngay cả khi không có giả thiết nào về các thành phần bất khả quy của các thớ
$X_s$, lược đồ $\operatorname{Pic}^{\tau}_{X/S}$ (sẽ được giới thiệu dưới
đây) vẫn có thể tồn tại.
\src{222}
\noindent b. (Murre). Cho $X$ là một lược đồ riêng trên một trường; khi đó
$\operatorname{Pic}_{X/k}$ tồn tại. Chứng minh dựa một phần trên chứng minh
của Chevalley [1] và sử dụng cốt yếu cấu trúc nhóm của hàm tử Picard.

Đối với một số bổ sung về lý thuyết các lược đồ Picard, đặc biệt trong mối
liên hệ với các lược đồ abel, nên tham khảo [7]. Cuối cùng, một thiếu sót đáng
kể của thuyết trình này là chưa có các «tiêu chuẩn tương đương» cho phép so
sánh lược đồ Picard của một lược đồ xạ ảnh với các lược đồ Picard của những
tiết diện siêu phẳng của nó; các định lý then chốt để xây dựng những tiêu
chuẩn như vậy nằm trong [5], và cần ghép thêm các định lý tồn tại của các
lược đồ Picard.

\fsection{1}{Các tính chất tôpô của các tiền lược đồ nhóm giao hoán}

Trước hết, cho $k$ là một trường và $G$ là một tiền lược đồ nhóm trên $k$.
Vì phần tử đơn vị $e$, do hữu tỷ trên $k$, nhất thiết là đóng, nên đường chéo
của $G\times_kG$ lập tức là đóng, và do đó $G$ là tách: mọi tiền lược đồ nhóm
trên một trường đều tách. Ta ký hiệu bởi $G^\circ$ thành phần liên thông của
phần tử đơn vị $e$. Vì $e$ hữu tỷ trên $k$, $G^\circ$ thực ra liên thông hình
học, nghĩa là phép tạo $G^\circ$ tương thích với mọi đổi trường cơ sở. Cũng
suy ra rằng $G^\circ$ ổn định dưới phép nhân khi xét như một tập; và khi $G$
Noether địa phương, nên $G^\circ$ mở, ta có thể xem $G^\circ$ như một nhóm
con mở của $G$. Sau đây ta giả sử $G$ thuộc kiểu hữu hạn địa phương trên $k$;
khi đó $G^\circ$ bất khả quy hình học và thuộc kiểu hữu hạn trên $k$. Thật
vậy, có thể giả sử $k$ đóng đại số, rồi giả sử thêm $G$ rút gọn (vì khi ấy
$G_{\mathrm{réd}}$ là một nhóm con của $G$, do
$G_{\mathrm{réd}}\times_kG_{\mathrm{réd}}$ vẫn rút gọn); do đó trơn trên
$k$ tại một tập mở không rỗng, và vì tịnh tiến tập mở ấy nên trơn ở mọi nơi.
Nhưng khi đó $G$ bất khả quy địa phương, nên các thành phần bất khả quy của nó
trùng với các thành phần liên thông; do đó $G^\circ$ bất khả quy. Cho $U$ là
một lân cận afin của $e$ trong $G^\circ$; dùng tính bất khả quy của $G^\circ$,
ta thấy ngay
$U\cdot U=G^\circ$, điều này chứng minh $G^\circ$ tựa compact, nên thuộc kiểu
hữu hạn trên $k$.

Để đơn giản, giả sử $G$ giao hoán. Với mỗi số nguyên $n>0$, cho $G^{(n)}$ là
ảnh ngược của $G^\circ$ qua đồng cấu $\varphi_n$ lũy thừa bậc $n$ trong
$G$; do đó $G^{(n)}$ là một nhóm con mở của $G$. Ta đặt:
\[
G^\tau=\bigcup_nG^{(n)}.
\]
\src{223}
\[
G^\sigma=\bigcup_{(n,p)=1}G^{(n)},
\qquad
G^\rho=\bigcup_hG^{(p^h)},
\]
trong đó $p$ là số mũ đặc trưng của trường $k$. Như vậy ta thu được các nhóm
con mở của $G$, thỏa mãn
\[
G^\sigma\cap G^\rho=G^\circ,
\qquad
G^\sigma G^\rho=G^\tau.
\]

\result{NHẬN XÉT}{} Ta có thể dựng lược đồ nhóm thương
$G/G^\circ=N$ (xem [3], IV), rồi định nghĩa $G^\tau$, $G^\sigma$, $G^\rho$
là các ảnh ngược trong $G$, tương ứng, của nhóm con xoắn của $N$, của tổng
các thành phần $q$-nguyên sơ với $q$ là số nguyên tố $\neq p$, và của thành
phần $p$-nguyên sơ. Cần lưu ý rằng $N$ là một lược đồ nhóm rời rạc khả ly trên
$k$; do đó, sau khi chọn một bao đóng đại số $\overline k$ của $k$, sinh ra
một nhóm Galois $\pi$, nó được đồng nhất với một nhóm rời rạc thông thường mà
$\pi$ tác động bằng các tự đẳng cấu. Chính bằng cách này ta có thể diễn giải
hiển nhiên việc dựng nhóm con xoắn và sự phân tích nhóm ấy thành các thành
phần $q$-nguyên sơ. Khi $G$ là lược đồ Picard của một lược đồ $X$ riêng trên
$k$, có thể gọi $N$ là \emph{lược đồ Néron--Severi} (rút gọn) của $X$ trên
$k$. Khi $G^\circ_{\mathrm{réd}}$ là một lược đồ nhóm con của $G^\circ$, điều
này đặc biệt xảy ra mỗi khi $k$ hoàn hảo hoặc $G^\circ$ riêng trên $k$ (chẳng
hạn khi $X$ chuẩn hình học), cũng nên đưa vào thương
$N'=G/G^\circ_{\mathrm{réd}}$; thương này có khuynh hướng ứng xử tốt hơn $N$
về phương diện chuyên biệt hóa, nghĩa là khi $X$ biến thiên trong một họ các
lược đồ đại số.

Bây giờ cho $S$ là một tiền lược đồ Noether địa phương và $G$ là một tiền
lược đồ nhóm trên $S$, thuộc kiểu hữu hạn địa phương trên $S$. Ta không giả
sử $G$ thuộc kiểu hữu hạn trên $S$, cũng không giả sử là tách trên $S$. Ta
đặt
\[
G^\circ=\bigcup_{s\in S}(G_s)^\circ,
\]
và, khi $G$ giao hoán,
\[
G^\tau=\bigcup_{s\in S}(G_s)^\tau,
\qquad
G^\sigma=\bigcup_{s\in S}(G_s)^\sigma,
\qquad
G^\rho=\bigcup_{s\in S}(G_s)^\rho.
\]
\src{224}
Đó là các tập con của $G$, ổn định dưới phép nhân của $G$; điều này hiển nhiên
không kéo theo rằng chúng có thể được định nghĩa bằng các tiền lược đồ nhóm
con của $G$. Đặc biệt, dường như nói chung không tồn tại tiền lược đồ nhóm
con nào của $G$ có tập nền là $G^\circ$. Dĩ nhiên, nếu một trong các tập ấy
mở thì, với cấu trúc cảm sinh, nó là một tiền lược đồ nhóm con mở của $G$. Ta
sẽ thấy điều này luôn đúng với $G^\tau$; vì thế, từ quan điểm các hàm tử biểu
diễn được, đặc biệt từ quan điểm «chuyên biệt hóa», tương đương số ứng xử thỏa
đáng hơn tương đương đại số. Sau đây là các tính chất tổng quát chính của
những tập vừa được định nghĩa:

\result{ĐỊNH LÝ}{1.1} $G^\circ$, $G^\tau$, $G^\sigma$, $G^\rho$ đều kiến
thiết được địa phương. Hơn nữa:

\noindent (i) $G^\circ$ tựa compact trên $S$. Khi các $G_s^\circ$ riêng và
$G$ tách trên $S$, thì $G^\circ$ riêng trên $S$, do đó đóng trong $G$.

\noindent (ii) $G^\tau$ mở. Nếu $G^\circ$ đóng thì $G^\tau$ cũng đóng.

\noindent (iii) Nếu $G^\circ$ đóng thì $G^\sigma$ cũng đóng, miễn là ta ở
trong đồng đặc trưng, nghĩa là mọi trường thặng dư của $S$ có cùng đặc trưng.
Nếu $G^\circ$ đóng và $G\longrightarrow S$ mở phổ dụng tại các điểm của
$G^\sigma$ (xem Hệ quả 1.5 dưới đây), thì
$G^\sigma\longrightarrow S$ mở phổ dụng.

\noindent (iv) Nếu $G^\circ$ đóng thì $G^\rho$ cũng đóng. Giả sử ta ở trong
đồng đặc trưng và, với mọi số nguyên $n>0$ sao cho $(n,p)=1$, đồng cấu lũy
thừa bậc $n$ trong $G$ là mở; khi đó $G^\rho$ mở.

Ta nêu vài chỉ dẫn cho chứng minh. Việc $G^\circ$ kiến thiết được địa phương
và tựa compact trên $S$ nằm trong bổ đề sau:

\result{BỔ ĐỀ}{1.2} Giả sử $S\neq\varnothing$. Khi đó tồn tại một tập mở
không rỗng $U$ trong $S$, một lược đồ nhóm $H$ thuộc kiểu hữu hạn trên $U$ có
các thớ liên thông, và một đồng cấu lược đồ nhóm
$H\longrightarrow G|U$ là một phép nhúng mở có ảnh $G^\circ|U$.

Để chứng minh bổ đề, có thể giả sử $S$ bất khả quy; gọi $\eta$ là điểm tổng
quát của nó và thực hiện đổi cơ sở
$S'=\operatorname{Spec}(\mathcal O_{S,\eta})\longrightarrow S$. Ta thu được
một lược đồ nhóm $G'$ trên vành Artin địa phương
$\mathcal O_{S,\eta}=A$, trong đó có một lược đồ nhóm con mở
$G'{}^\circ$ thuộc kiểu hữu hạn trên $A$, như đã nói ở trên. Vì vậy nó thu
được từ một lược đồ nhóm $H$ thuộc kiểu hữu hạn trên một lân cận mở $U$ của
$\eta$, và phép nhúng chính tắc $G'{}^\circ\longrightarrow G'$ thu được từ
một phép nhúng
\src{225}
mở $H\longrightarrow G|U$, và phép nhúng này sẽ là một đồng cấu lược đồ nhóm
khi $U$ đủ nhỏ. Các thớ của $H$ liên thông nếu lấy $U$ đủ nhỏ và đều có
cùng chiều, bằng chiều của các thớ của $G$ khi $U$ đủ nhỏ; do đó, với mọi
$s\in U$, ảnh của
$H_s$ trong $G_s$ chính là $G_s^\circ$ (sau khi lại thu nhỏ $U$), chứng minh
1.2.

Khẳng định thứ hai trong (i) nằm trong bổ đề sau (áp dụng cho một lân cận mở
tựa compact của $G^\circ$ trong $G$):

\result{BỔ ĐỀ}{1.3} Cho $X$ là một tiền lược đồ thuộc kiểu hữu hạn và tách
trên $S$ Noether địa phương, $g$ là một tiết diện của $X$ trên $S$, và
$X^\circ$ là hợp của các thành phần liên thông của $g(s)$ trong các $X_s$.
Cho $s\in S$ sao cho $X_s^\circ$ riêng trên $k(s)$. Khi đó tồn tại một lân
cận mở $U$ của $s$ sao cho $X^\circ|U$ riêng trên $U$, và do đó đóng trong
$X|U$.

Bằng hạ phẳng trung thành trên cơ sở, ta quy về trường hợp $S$ là phổ của một
vành địa phương đầy đủ và $s$ là điểm đóng của nó. Áp dụng [2], III, 5.5.1,
ta thấy $X$ phân tích thành tổng của hai tập mở rời nhau $X'$ và $X''$, trong
đó tập thứ nhất riêng trên $S$ và thỏa $X'_s=X_s^\circ$. Điều này quy ta về
trường hợp $X=X'$, nghĩa là $X$ riêng trên $S$. Trong trường hợp ấy, ta dùng
một chứng minh chuẩn dựa trên tiêu chuẩn định giá cho tính riêng của một bộ
phận (bị bỏ sót trong [2], chương II).

Ta chứng minh $G^\tau$ mở; hay tương đương, vì phép tạo $G^\tau$ (cũng như
phép tạo $G^\circ$, $G^\sigma$, $G^\rho$) giao hoán với mở rộng cơ sở: với
mọi tiết diện $g$ của $G$ trên $S$, $g^{-1}(G^\tau)$ mở. Điều này có nghĩa là
hai điều:

\medskip
\noindent a. Cho $y\in S$ sao cho $g(y)\in G_y^\tau$; khi đó, với mọi
$y'\in\overline{\{y\}}$ đủ gần $y$, ta có $g(y')\in G_{y'}^\tau$.

\medskip
\noindent b. Cho $y'\in S$ sao cho $g(y')\in G_{y'}^\tau$; khi đó, với mọi
tổng quát hóa $y$ của $y'$, ta có $g(y)\in G_y^\tau$.

Đối với (a), lưu ý rằng tồn tại một số nguyên $n>0$ sao cho
$g^n(y)\in G_y^\circ$; vì $G^\circ$ kiến thiết được,
$(g^n)^{-1}(G^\circ)$ kiến thiết được, nên
$g^n(y')\in G_{y'}^\circ$ với $y'\in\overline{\{y\}}$ đủ gần $y$. Đối với
(b), lưu ý rằng các $g(y')^n=g^n(y')$ nằm trong một tập mở tựa compact của
$G_{y'}$ (vì chúng thuộc một số hữu hạn lớp modulo $G_{y'}^\circ$); do đó tồn
tại một tập mở tựa compact $U$ trong $G$ chứa các $g^n(y')$, và vì thế cũng
chứa các tổng quát hóa $g^n(y)$ của chúng. Vậy các lũy thừa của $g(y)$ nằm
trong một tập mở tựa compact của $G_y$, điều này dễ dàng kéo theo
$g(y)\in G_y^\tau$.
\src{226}
Giả sử $G^\circ$ đóng; ta chứng minh $G^\tau$ cũng đóng. Vì đã biết
$G^\tau$ mở, nên kiến thiết được địa phương, chỉ còn phải chứng minh nó ổn
định dưới phép chuyên biệt hóa. Điều này suy ra từ việc nó là hợp của các tập
đóng, cụ thể là các ảnh ngược, qua các đồng cấu $\varphi_n$ lũy thừa bậc
$n$, của tập đóng $G^\circ$.

Cũng lập luận ấy sẽ chứng minh $G^\sigma$ và $G^\rho$ đóng nếu $G^\circ$
đóng (trong trường hợp thứ nhất, với điều kiện ta ở trong đồng đặc trưng), một
khi đã chứng minh $G^\sigma$ và $G^\rho$ kiến thiết được địa phương. Thật vậy,
với $x\in G^\tau$, gọi $\nu(x)$ là số nguyên nhỏ nhất $n>0$ sao cho đồng cấu
$\varphi_n$ lũy thừa bậc $n$ đưa $x$ vào $G^\circ$. Khi đó $G^\sigma$
(tương ứng, $G^\rho$) gồm các $x\in G$\footnote{Bản in ghi «$X$» ở đây,
nhưng ký hiệu ấy không được định nghĩa trong đoạn này; cả bốn locus đang xét đều là
các bộ phận của $G$.} sao cho $\nu(x)$ nguyên tố cùng nhau với $p$ (tương ứng,
là một lũy thừa của $p$), và khẳng định của ta về tính kiến thiết được suy ra
từ khẳng định chính xác hơn sau đây:

\result{BỔ ĐỀ}{1.4} Hàm $\nu$ trên $G^\tau$ kiến thiết được địa phương.

Thật vậy, điều này có nghĩa là, với mọi số nguyên $n>0$, tập các
$x\in G^\tau$ sao cho $\nu(x)=n$ kiến thiết được địa phương; nhưng đó là hiệu
của $\varphi_n^{-1}(G^\circ)$ và hợp của các
$\varphi_d^{-1}(G^\circ)$, trong đó $d$ chạy qua các ước thực sự của $n$. Vì
$G^\circ$ kiến thiết được địa phương, các $\varphi_d^{-1}(G^\circ)$ cũng vậy,
và do đó hiệu nói trên cũng vậy.

Giả sử $G^\circ$ đóng và $G\longrightarrow S$ mở phổ dụng tại các điểm của
$G^\sigma$; ta chứng minh $G^\sigma\longrightarrow S$ là mở, nghĩa là biến
một lân cận trong $G^\sigma$ của $x\in G^\sigma$ thành một lân cận của
$y=f(x)$. Vì $G^\sigma$ kiến thiết được địa phương, ta quy về việc chứng minh
rằng, với mọi tổng quát hóa $y'$ của $y$, tồn tại một tổng quát hóa $x'$ của
$x$ trong $G^\sigma$ nằm trên $y'$. Bằng đổi cơ sở, điều này quy ta về trường
hợp $S$ là phổ của một vành định giá rời rạc, còn $y,y'$ lần lượt là điểm đóng
và điểm tổng quát. Dùng việc $G\longrightarrow S$ mở tại $x$ (do đó tồn tại
một tổng quát hóa $x_1$ của $x$ trong $G$ nằm trên $y'$), ta có thể giả sử
thêm rằng tồn tại một tiết diện $g$ của $G$ trên $S$ sao cho $x=g(y)$ (sau
khi lại đổi cơ sở nếu cần). Nếu $k(y')$ có đặc trưng 0, chỉ cần lấy một tổng
quát hóa bất kỳ $x'$ của $x$ trong $G$ nằm trên $y'$; nó thuộc $G^\tau$ vì
$G^\tau$ mở, và do đó thuộc $G^\sigma$ vì
$G_{y'}^\sigma=G_{y'}^\tau$. Nếu $k(y')$ có đặc trưng $p>0$, đặt
\[
\nu(g(y'))=p^hm\qquad\text{với }(m,p)=1 ;
\]
\src{227}
chọn hai số nguyên $a$ và $b$ sao cho $ap^h+bm=1$, và đặt
\[
g_1=g^{ap^h},\qquad g_2=g^{bm},\qquad\text{do đó}\qquad g=g_1g_2 ;
\]
theo cách dựng, ta có $g_1(y')\in G^\sigma$ và
$g_2(y')\in G^\rho$. Vì $G^\circ$ đóng, suy ra
$g_1(y)\in G^\sigma$ và $g_2(y)\in G^\rho$; mặt khác, từ
\[
g(y)=g_1(y)g_2(y)\in G^\sigma,
\]
ta cũng có $g_2(y)\in G^\sigma$, nên
\[
g_2(y)\in G_y^\sigma\cap G_y^\rho=G_y^\circ.
\]
Từ giả thiết và việc $S$ là phổ của một vành định giá rời rạc, suy ra
$G-(G_y-G_y^\circ)$ là một tập mở của $G$ mà trên đó
$G\longrightarrow S$ cảm sinh một cấu xạ mở; do đó, qua mỗi điểm của
$G_y^\circ$ có một «tựa tiết diện». Vì vậy, sau khi lại thực hiện một mở rộng
trên cơ sở $S$ nếu cần, ta có thể giả sử tồn tại một tiết diện $g'_2$ của
$G^\circ$ trên $S$ sao cho $g'_2(y)=g_2(y)$. Đặt $g'=g_1g'_2$; theo cách dựng,
$g'(y)=g(y)=x$, và
\[
g'(y')=g_1(y')g'_2(y')\in G^\sigma,
\]
do đó $g'(y')=x'$\footnote{Bản in ghi «$g(y')=x'$» ở đây; tiết diện được dựng
trong câu trước và đang được tính thớ tổng quát là $g'$.} là một tổng quát hóa
của $x$ trong $G^\sigma$ nằm trên $y'$, chứng minh (iii).

Cuối cùng, ta chứng minh khẳng định cuối của (iv). Ta quy về việc chứng minh
rằng, nếu $x\in G^\rho$, thì mọi tổng quát hóa $x'$ của $x$ đều thuộc
$G^\rho$. Có thể giả sử (sau khi lấy ảnh qua $\varphi_{p^h}$ với $h$ thích
hợp nếu cần)\footnote{Bản in ghi $\varphi_{n^h}$; đối với $x\in G^\rho$, định
nghĩa ở tr.~223 cho một lũy thừa của số mũ đặc trưng $p$ đưa $x$ vào
$G^\circ$.} rằng thậm chí $x\in G^\circ$. Khi đó, với mọi số nguyên $n$
nguyên tố cùng nhau với đặc trưng, $x$ thuộc ảnh của $\varphi_n$ (vì lũy thừa
bậc $n$ trong một nhóm liên thông thuộc kiểu hữu hạn trên một trường có đặc
trưng nguyên tố cùng nhau với $n$ là toàn ánh). Vì $\varphi_n$ mở, suy ra $x'$
cũng thuộc ảnh của $\varphi_n$. Chính xác hơn, cho $U$ là một tập mở tựa
compact của $G^\tau$ chứa $G_y^\circ$; khi đó
$x'\in\varphi_n(U)$ với mọi $n$ nguyên tố cùng nhau với $p$. Lấy $n$ là một
bội chung của các thừa số nguyên tố cùng nhau với $p$ của $\nu(z)$, với
$z\in U$, ta thu được $x'\in G^\rho$.

Khẳng định (1.1), (iii), được bổ sung như sau:

\result{HỆ QUẢ}{1.5} Cho $n>1$ là một số nguyên sao cho đồng cấu lũy thừa bậc
$n$, $\varphi_n:G\longrightarrow G$, mở phổ dụng (chẳng hạn étale), và đặt
\[
G^{(n)}=\varphi_n^{-1}(G^\circ),
\]
Giả sử trong mỗi thớ hình học $G_{\bar s}$, tập con
\src{228}
$\bigcup_{h\geq0}|{}_{n^h}G_{\bar s}|$ trù mật trong
$\bigcup_{h\geq0}G_{\bar s}^{(n^h)}$. Khi đó
$G\longrightarrow S$ mở phổ dụng tại các điểm của các $G^{(n^h)}$. Đặc biệt,
nếu các $G_s^\circ$ không chứa thành phần cộng tính và, với mọi số nguyên
$n>1$, đồng cấu $\varphi_n$ mở phổ dụng tại các điểm có đặc trưng thặng dư
nguyên tố cùng nhau với $n$\footnote{Bản in ghi «nguyên tố cùng nhau với
$p$». Trong mệnh đề được lượng hóa theo $n$ này, tiêu chuẩn étale và lập luận
bằng các hạt nhân ${}_nG$ đòi hỏi đặc trưng thặng dư nguyên tố cùng nhau với
$n$.}, thì $G\longrightarrow S$ mở phổ dụng tại các điểm của $G^\sigma$, và
do đó, theo (1.1), (iii), nếu $G^\circ$ đóng thì
$G^\sigma\longrightarrow S$ mở phổ dụng. Trong các trường hợp này,
$s\longmapsto\dim G_s$ là một hàm hằng địa phương trên $S$.

Thật vậy, từ giả thiết suy ra hạt nhân ${}_nG$ của $\varphi_n$ mở phổ dụng
trên $S$, nên $G\longrightarrow S$ mở phổ dụng tại các điểm của ${}_nG$, và
do đó (thay $n$ bởi $n^h$) tại các điểm của ${}_{n^h}G$. Giả thiết về tính trù mật nói chính xác rằng, trên từng thớ, các điểm có cấp
là một lũy thừa của $n$ trù mật trong hợp của các $G^{(n^h)}$. Đối với khẳng
định riêng, việc không có thành phần cộng tính kéo theo tính trù mật này cho
các số nguyên $n$ khả nghịch trong trường thặng dư; từ đó dễ suy
ra $G\longrightarrow S$ mở phổ dụng tại mọi điểm của hợp này, đặc biệt dọc
theo tiết diện đơn vị, và rồi dễ suy ra $s\longmapsto\dim G_s$ là một hàm
hằng địa phương.

\result{NHẬN XÉT}{1.6} Nhắc lại rằng một cấu xạ $X\longrightarrow S$ được gọi
là mở phổ dụng nếu nó biến mọi tập mở thành một tập mở và vẫn giữ tính chất
đáng khen ấy sau mọi phép đổi cơ sở $S'\longrightarrow S$. Điều này cũng có
nghĩa (khi $S$ Noether địa phương và $X\longrightarrow S$ thuộc kiểu hữu hạn
địa phương) rằng mọi thành phần bất khả quy của $X$ trội lên $S$, và tính chất
này được bảo toàn sau mọi phép đổi cơ sở $S'\longrightarrow S$. Trong cả hai
khẳng định ấy, hơn nữa, chỉ cần (với các giả thiết hữu hạn ở trên) kiểm tra
bằng các phép đổi cơ sở $S'\longrightarrow S$, trong đó $S'$ là phổ của một
vành định giá rời rạc (đầy đủ với trường thặng dư đóng đại số nếu muốn~...).
Định nghĩa mở rộng hiển nhiên cho trường hợp một bộ phận $Z$ của $X$ (chẳng
hạn bộ phận $G^\sigma$ của $G$). Hoàn toàn bình thường là, ngay cả khi bắt đầu
từ một cấu xạ xạ ảnh trơn $X\longrightarrow S$ có các thớ hình học liên thông
(chẳng hạn tích thớ của họ môđun các đường cong elliptic trên
$S=\operatorname{Spec}\mathbf C[j]$ với chính nó),
$\operatorname{Pic}_{X/S}$ vẫn không mở phổ dụng trên $S$, nghĩa là có thể có
những thành phần bất khả quy của $\operatorname{Pic}_{X/S}$ nằm hoàn toàn
trên một điểm duy nhất của $S$: điều này liên quan đến việc hạng của nhóm
Néron--Severi của các thớ của $X/S$ có thể nhảy lên (các hiện tượng «nhân
phức»). Trái lại, (1.5) bảo đảm rằng, trong những trường hợp thuận lợi,
$\operatorname{Pic}_{X/S}^\sigma$ (và dường như thường ngay cả
$\operatorname{Pic}_{X/S}^\tau$) mở phổ dụng trên $S$.
\src{229}
Cuối cùng, sau đây là một trường hợp hữu ích trong đó, một cách ngoại lệ,
$G^\circ$ lại là mở:

\result{HỆ QUẢ}{1.7} Giả sử $G\longrightarrow S$ mở phổ dụng tại các điểm của
$G^\circ$ (xem (1.5)), và các thớ $G_s$ khả ly, do đó trơn trên $k(s)$ (điều
kiện sau tự động được thỏa mãn khi đặc trưng thặng dư bằng không, theo một kết
quả của Cartier). Khi đó $G^\circ$ mở trong $G$. Nếu thêm nữa $S$ rút gọn, thì
$G^\circ$ trơn (đặc biệt là phẳng) trên $S$.

Khẳng định thứ nhất còn có thể được làm chính xác hơn như sau. Nếu, với một
$s\in S$ cho trước, $G\longrightarrow S$ mở phổ dụng tại các điểm của
$G_s^\circ$ và $G_s^\circ$ khả ly trên $k(s)$, thì $G^\circ$ là một lân cận
của $G_s^\circ$ (hơn nữa, giả thiết đặt tại $s$ vẫn sẽ đúng tại các điểm lân
cận). Mệnh đề này, vốn độc lập với mọi cấu trúc nhóm, sẽ được chứng minh trong
[2], IV, đoạn 7. Khẳng định cuối của (1.7), cũng độc lập với mọi cấu trúc nhóm
và mọi vấn đề về thành phần liên thông, là một trường hợp riêng của một tiêu
chuẩn phẳng được cho trong [2], IV, đoạn 5, và tiêu chuẩn này tổng quát hơn còn
kéo theo:

\result{HỆ QUẢ}{1.8} Cho $U$ là một bộ phận mở của một thớ $G_s$, sao cho
$G\longrightarrow S$ mở phổ dụng tại các điểm của $U$ (xem (1.5)). Nếu $G_s$
khả ly trên $k(s)$ và $S$ rút gọn tại $s$, thì $G$ phẳng trên $S$ tại các điểm
của $U$ (do đó, trong trường hợp này, trơn trên $S$ tại các điểm của $U$).

\result{CÁC NHẬN XÉT}{1.9} Tôi không biết, khi $G$ tách trên $S$, liệu
$G^\circ$ hay $G^\tau$ có luôn đóng trong $G$ hay không; điều này có vẻ khó xảy
ra. Dù sao, ta có những phản ví dụ hiển nhiên nếu bỏ giả thiết tách, chẳng hạn
lấy đường thẳng afin mà gốc được tách đôi vô hạn đếm được lần. Như vậy ta thu
được một lược đồ nhóm $G$ trên đường thẳng afin, mọi thớ của nó, trừ một thớ
bằng $\mathbf Z$, đều rút về nhóm chỉ gồm phần tử đơn vị. Lược đồ nhóm này là
con mở, cụ thể là bao đóng của tiết diện đơn vị, trong tiền lược đồ Picard của
$S$-lược đồ $X$ tương ứng với một họ conic suy biến thành hai đường thẳng giao
nhau.

Hơn nữa, ngay cả khi bắt đầu với một lược đồ nhóm hữu hạn và phẳng $G$ trên
phổ $S$ của một vành định giá rời rạc $V$, sao cho $G^\circ$ chỉ còn là tiết
diện đơn vị và do đó là đóng, nhiều kết luận vẫn trở thành sai nếu bỏ một số
giả thiết. Giả sử $V$ có đồng đặc trưng $p>0$, và cho $G$ là hạt nhân của đồng
cấu
$\mathbf G_a\longrightarrow\mathbf G_a$ được xác định bởi đồng cấu các hàm tử
\[
f\longmapsto f^p-tf,
\]
trong đó $t$ là một phần tử đều hóa của $V$.
\src{230}
(Nhắc lại rằng, theo định nghĩa, «nhóm cộng» $\mathbf G_a$ trên $S$ biểu diễn
hàm tử $S'\longmapsto\Gamma(S',\mathcal O_{S'})$.) Khi đó, với tư cách là một
$S$-lược đồ,
\[
G=\operatorname{Spec}V[x]/(x^p-tx)
\]
là hợp của $p$ «đường thẳng» giao nhau, mô tả một nhóm
$\mathbf Z/p\mathbf Z$ suy biến thành một nhóm vô cùng bé kiểu cộng. Trong ví
dụ này, $G^\circ=G^\sigma$ (và đó là một trong $p$ đường thẳng), nên không mở
trong $G$. Trong trường hợp khác đặc trưng, với đặc trưng thặng dư $p>0$, ta
bắt đầu từ lược đồ nhóm $H=\symbf{\mu}_p$, là hạt nhân của phép lấy lũy thừa
bậc $p$ trong $\mathbf G_m$, do đó
\[
H=\operatorname{Spec}V[x]/(x^p-1) ;
\]
đây lại là hợp của $p$ đường thẳng giao nhau, mô tả một nhóm
$\mathbf Z/p\mathbf Z$ (trong đặc trưng 0) suy biến thành một nhóm vô cùng bé
kiểu nhân trong đặc trưng $p$. Cho $H'$ là «lược đồ nhóm hằng» được xác định
bởi nhóm hữu hạn thông thường $\mathbf Z/p\mathbf Z$, tức hợp rời của
$\mathbf Z/p\mathbf Z$ bản sao của $S$, hay cũng là
$\operatorname{Spec}V^{\mathbf Z/p\mathbf Z}$. Khi đó $G=H\times_S H'$ mô tả
một nhóm $(\mathbf Z/p\mathbf Z)^2$ trong đặc trưng 0, suy biến thành tích của
một nhóm vô cùng bé với $\mathbf Z/p\mathbf Z$ trong đặc trưng $p$. Ở đây
$G^\rho$ là hợp của tiết diện đơn vị và thớ đặc biệt, nên $G^\rho$ không mở,
trong khi $G^\sigma$ là hợp của tiết diện đơn vị và thớ tổng quát, nên không
đóng, trái với điều được khẳng định trong các trường hợp đồng đặc trưng ở
(1.1), (iii) và (iv). Dĩ nhiên, đây là những hiện tượng gắn với đặc trưng
$p>0$. Thật vậy, các kết quả trên cho:

\result{HỆ QUẢ}{1.10} Giả sử mọi đặc trưng thặng dư của $S$ đều bằng không,
do đó $G^\tau=G^\sigma$ và $G^\circ=G^\rho$. Khi đó $G^\tau=G^\sigma$ là mở,
thậm chí vừa mở vừa đóng nếu $G$ tách trên $S$ và các $G_s^\circ$ riêng; dưới
cùng giả thiết ấy, $G^\circ=G^\rho$ riêng trên $S$, nên đóng trong $G$. Cuối
cùng, giả sử với mọi số nguyên $n>1$, đồng cấu lũy thừa bậc $n$ trong $G$ mở
phổ dụng; khi đó $G^\circ$ mở. Nếu thêm nữa các $G_s^\circ$ không có thành
phần cộng (chẳng hạn các $G_s^\circ$ riêng, như trên), thì
$G^\tau\longrightarrow S$ mở phổ dụng, và thậm chí trơn nếu $S$ rút gọn.

Cuối cùng, ta ghi nhận kết quả dễ sau đây:

\result{MỆNH ĐỀ}{1.11} Tồn tại một bộ phận mở $U$ của $S$ sao cho tập $G'$ gồm
các điểm của $G$ tại đó $G$ trơn (tương ứng phẳng) trên $S$ là tập nền của một
lược đồ nhóm con mở cảm sinh từ $G|U$. Hơn nữa, mọi tiết diện của $G$ trên $U$
đều là một tiết diện của $G'$ trên $U$.
\src{231}
\result{HỆ QUẢ}{1.12} Nếu $G$ trơn (tương ứng phẳng) trên $S$ tại các điểm của
tiết diện đơn vị, thì nó cũng như vậy tại các điểm của mọi tiết diện của $G$
trên $S$, và tại mọi điểm của $G^\circ$. Nếu thêm nữa, với mọi số nguyên
$n>0$, đồng cấu lũy thừa bậc $n$
$\varphi_n:G\longrightarrow G$ là étale tại các điểm có đặc trưng nguyên tố
cùng nhau với $n$, thì $G$ trơn (tương ứng phẳng) trên $S$ tại mọi điểm của
$G^\sigma$.

\fsection{2}{Ứng dụng vào các tính chất địa phương của các lược đồ Picard}

\result{ĐỊNH LÝ}{2.1}

\noindent(i) Cho $f:X\longrightarrow S$ là một cấu xạ riêng và trơn, và giả
sử $\operatorname{Pic}_{X/S}$ tồn tại (chẳng hạn $f$ là một cấu xạ xạ ảnh).
Khi đó $\operatorname{Pic}_{X/S}$ tách trên $S$, và với mọi bộ phận đóng $Z$
của $\operatorname{Pic}_{X/S}$ thuộc kiểu hữu hạn trên $S$, $Z$ riêng trên
$S$.

\noindent(ii) Cho $X$ là một tiền lược đồ trên một trường $k$, riêng và chuẩn
hình học. Khi đó $\operatorname{Pic}_{X/k}^\circ$\footnote{Bản in ghi
$\operatorname{Pic}_{X/S}^\circ$, mặc dù $S$ không xuất hiện trong mệnh đề
trên trường $k$ này.} riêng trên $k$.

\begin{proof}
(i) Các tiêu chuẩn định giá ([2], II, đoạn 7) quy ta về việc chứng minh điều
sau: nếu $S$ là phổ của một vành định giá rời rạc đầy đủ, $U$ là tập mở chỉ
gồm điểm tổng quát của $S$, thì mọi tiết diện hữu tỷ của
$\operatorname{Pic}_{X/S}$ trên $S$, nghĩa là mọi tiết diện trên $U$, đều kéo
dài duy nhất thành một tiết diện trên $S$. Theo định nghĩa của
$\operatorname{Pic}_{X/S}$, điều này tương đương với mệnh đề sau: với mọi
Môđun khả nghịch $\mathcal L$ trên $V=f^{-1}(U)$, tồn tại một Môđun khả nghịch
trên $X$ kéo dài $\mathcal L$, duy nhất sai khác một đẳng cấu. Điều này dễ dàng
suy ra từ mô tả các Môđun khả nghịch trên $V$, tương ứng trên $X$, bằng các
lớp «ước Cartier», vì các vành địa phương của $X$ là chính quy ($X$ trơn trên
lược đồ chính quy $S$), nên là các vành phân tích duy nhất theo Auslander; do
đó mọi ước trên $X$\footnote{Bản in ghi ở đây «trên $S$»; tính phân tích duy
nhất được viện dẫn là của các vành địa phương của $X$.} đều là một ước
Cartier. Thật vậy, mọi ước trên $V$ đều có thể được kéo dài thành một ước trên
$X$ bằng cách lấy «bao đóng» của nó.
\end{proof}

\result{NHẬN XÉT}{2.2} Chứng minh vẫn đúng nếu chỉ giả sử $f$ phẳng và các thớ
$X_s$ là các giao đầy đủ địa phương và trơn trên $k(s)$ ở đối chiều $\leq2$,
theo kết quả sau được chứng minh trong [5]: một vành địa phương Noether là giao
đầy đủ và chính quy ở đối chiều $\leq3$ thì là vành phân tích duy nhất («phỏng
đoán Samuel»). Cần lưu ý rằng
\src{232}
kết quả trở thành sai nếu thay «đối chiều $\leq2$» bằng «đối chiều
$\leq1$», nghĩa là bằng giả thiết «chuẩn», như có thể thấy từ ví dụ một họ các
siêu mặt bậc hai không kỳ dị suy biến thành một nón bậc hai.

(ii) Dùng bổ đề Chow, ta quy về trường hợp $X$ xạ ảnh, do đó được nhúng vào
$\mathbf P_k^n$; ta có thể giả sử $X$ liên thông. Nếu $\dim X=1$, thì $X$ trơn
trên $k$, và ta áp dụng (i). Nếu $\dim X\geq2$, ta có thể dùng một biến thể của
các «tiêu chuẩn tương đương» đã biết, từ đó suy ra rằng tồn tại hữu hạn đường
cong $Y_i$ trơn nằm trong $X$ (thu được bằng cách giao $X$ với các không gian
con tuyến tính thích hợp của $\mathbf P_k^n$), sao cho
\[
\operatorname{Pic}_{X/k}^\tau\longrightarrow
\prod_i\operatorname{Pic}_{Y_i/k}^\tau
\]
là một đơn cấu, và do đó cảm sinh một đơn cấu trên các thành phần liên thông.
Vì thành phần liên thông của vế phải riêng trên $k$ theo kết quả trên, và vì
đồng cấu đang xét là một đồng cấu các lược đồ nhóm, tất yếu là một phép nhúng
đóng, nên $\operatorname{Pic}_{X/k}^\circ$ cũng riêng trên $k$. Có thể tránh
việc dùng các tiêu chuẩn tương đương tinh tế bằng cách sử dụng cấu trúc của các
nhóm đại số giao hoán trên một trường đóng đại số (do Chevalley--Borel); ta quy
về việc chứng minh rằng mọi cấu xạ từ đường thẳng afin khuyết
gốc\footnote{Bản in ghi «từ đường thẳng afin». Bản đính chính liên kết,
tr.~303, quy định «đường thẳng afin khuyết gốc».} vào
$\operatorname{Pic}_{X/k}^\tau$ là hằng. Điều này tương đương với việc mọi
Môđun khả nghịch trên $X[t,t^{-1}]$\footnote{Bản in ghi $X[t]$; bản đính
chính liên kết, tr.~303, quy định $X[t,t^{-1}]$.} đều đến từ một Môđun khả
nghịch trên $X$. Kết quả quen thuộc và rất sơ cấp này thậm chí không dùng đến
việc $X$ riêng trên $k$ (giả thiết $X$ chuẩn cho phép lập tức quy về trường hợp
$X$ chính quy).

\result{HỆ QUẢ}{2.3} Cho $f:X\longrightarrow S$ là một cấu xạ riêng và chuẩn
(nghĩa là phẳng với các thớ hình học chuẩn), và giả sử
$\operatorname{Pic}_{X/S}$ tồn tại. Khi đó $\operatorname{Pic}_{X/S}^\circ$
riêng trên $S$, nên đóng trong $\operatorname{Pic}_{X/S}$; hơn nữa,
$\operatorname{Pic}_{X/S}^\tau$ và $\operatorname{Pic}_{X/S}^\rho$ cũng đóng,
và $\operatorname{Pic}_{X/S}^\sigma$ cũng vậy trong «đồng đặc trưng».

Ta áp dụng (1.1) và (2.1), (ii).

\result{HỆ QUẢ}{2.4} Cho $f:X\longrightarrow S$ là một cấu xạ riêng và trơn,
sao cho $\operatorname{Pic}_{X/S}$ tồn tại và là hợp rời của các lược đồ
$P^{(i)}$ thuộc kiểu hữu hạn trên $S$ (xem [3], V, 4.1). Khi đó mỗi $P^{(i)}$
riêng trên $S$.

Điều này suy ra từ (2.1), (i). Như đã nêu trong (2.2), kết quả có thể được tổng
quát hóa bằng cách đặt các giả thiết ít hạn chế hơn lên các thớ của $f$, nhưng
trở thành sai nếu ta chỉ giả sử $f$ chuẩn. Trong trường hợp này, tôi cũng không
biết liệu $\operatorname{Pic}_{X/S}^\tau$ có vẫn riêng trên $S$ hay không,
ngay cả khi giả sử nó thuộc kiểu hữu hạn trên $S$.
\src{233}
\result{ĐỊNH LÝ}{2.5} Cho $f:X\longrightarrow S$ là một cấu xạ riêng và
phẳng, sao cho $\operatorname{Pic}_{X/S}$ tồn tại, và với mỗi số nguyên $n$,
cho $\varphi_n$ là đồng cấu lũy thừa bậc $n$ trong tiền lược đồ nhóm này. Khi
đó $\varphi_n$ là étale tại mọi điểm
$x\in\operatorname{Pic}_{X/S}$\footnote{Bản in ghi «$x\in X$»;
$\varphi_n$ là tự đồng cấu lũy thừa của tiền lược đồ Picard.} có đặc trưng
thặng dư nguyên tố cùng nhau với $n$.

Theo đặc trưng hóa vi phân của các cấu xạ étale, khẳng định này tương đương với
khẳng định sau:

\result{BỔ ĐỀ}{2.6} Giả sử $S$ là phổ của một vành Artin địa phương $A$ mà
iđêan cực đại $\mathfrak m$ có lũy thừa bậc $(\nu+1)$ bằng không, và đặt
\[
A_{\nu-1}=A/\mathfrak m^\nu,
\qquad X_{\nu-1}=X\otimes_A A_{\nu-1},
\]
$\mathcal L$ là một Môđun khả nghịch trên $X$, và cuối cùng
$\mathcal L'_{\nu-1}$ là một Môđun khả nghịch trên $X_{\nu-1}$ mà lũy thừa
tensor bậc $n$ đẳng cấu với
\[
\mathcal L_{\nu-1}=\mathcal L\otimes_A A_{\nu-1}.
\]
Khi đó tồn tại một Môđun khả nghịch $\mathcal L'$ trên $X$ mà lũy thừa tensor
bậc $n$ đẳng cấu với $\mathcal L$ (nếu $n$ nguyên tố cùng nhau với đặc trưng
thặng dư của $k=k(A)$).

Ta chứng minh bổ đề. Đặt
\[
V=\mathfrak m^\nu=\mathfrak m^\nu/\mathfrak m^{\nu+1} ;
\]
đây là một không gian vectơ trên $k=k(A)$. Trước hết ta tìm cách kéo dài
$\mathcal L'_{\nu-1}$ thành một Môđun khả nghịch tùy ý $\mathcal L'$ trên
$X$. Trở ngại cho việc này nằm trong
\[
H^2(X_0,\mathcal O_{X_0})\otimes_k V,
\]
nhưng theo giả thiết về $\mathcal L'_{\nu-1}$ và sự kiện
$\mathcal L_{\nu-1}$ có thể được kéo dài, ta thấy tích của trở ngại này với
$n$ bằng không; do đó chính nó bằng không vì $n$ nguyên tố cùng nhau với đặc
trưng. Hơn nữa, độ tùy ý trong phép kéo dài đã chọn nằm trong
\[
H^1(X_0,\mathcal O_{X_0})\otimes_k V,
\]
và sai lệch $\xi$ của $\mathcal L'^{\otimes n}$ so với $\mathcal L$ nằm trong
cùng Môđun ấy. Nếu ta tìm cách hiệu chỉnh $\mathcal L'$ để làm sai lệch này
bằng không, thì bài toán quy về việc tìm một $\eta$ trong Môđun nói trên sao
cho $n\eta=\xi$. Điều này lại khả thi vì $n$ nguyên tố cùng nhau với đặc
trưng.

\result{HỆ QUẢ}{} Dưới các điều kiện của (2.5), giả sử thêm rằng các lược đồ
Picard của các thớ $X_s$ không chứa thành phần cộng (chẳng hạn các $X_s$ chuẩn
hình học, xem (2.1), (ii)). Khi đó
$\operatorname{Pic}_{X/S}\longrightarrow S$ mở phổ dụng tại các điểm của
$\operatorname{Pic}_{X/S}^\circ$. Nếu $\operatorname{Pic}_{X/S}^\circ$ đóng
(chẳng hạn các $X_s$ chuẩn hình học, xem (2.3)), thì chính
$\operatorname{Pic}_{X/S}^\sigma$ mở phổ dụng trên $S$. Cuối cùng, trong
trường hợp đồng đặc trưng, $\operatorname{Pic}_{X/S}^\rho\longrightarrow S$
mở phổ dụng.

Ta áp dụng (1.5) và (1.1).
\src{234}
\result{HỆ QUẢ}{2.7} Cho $f:X\longrightarrow S$\footnote{Bản in ghi
«$f:X\longrightarrow Y$», mặc dù tiền lược đồ Picard và hàm số sau đều có cơ
sở $S$.} là một cấu xạ riêng và phẳng, sao cho $\operatorname{Pic}_{X/S}$ tồn
tại. Khi đó hàm số
\[
s\longmapsto\dim\operatorname{Pic}_{X_s/k(s)}
\]
trên $S$ là nửa liên tục trên (nghĩa là nó có thể nhảy lên, nhưng không thể
nhảy xuống), và thậm chí liên tục (nghĩa là hằng địa phương) nếu các
$\operatorname{Pic}_{X_s/k(s)}$ không chứa thành phần cộng.

Khẳng định thứ nhất đúng một cách tầm thường, hoặc gần như vậy, đối với mọi
tiền lược đồ nhóm thuộc kiểu hữu hạn địa phương trên một cơ sở Noether địa
phương, vì chỉ cần xét dọc theo tiết diện đơn vị. Khẳng định thứ hai suy ra từ
(2.5).

\result{NHẬN XÉT}{2.8} Cho $s,s'\in S$ sao cho $s$ là một chuyên biệt hóa của
$s'$. Khi đó (2.7) tương đương với một bất đẳng thức (tương ứng một đẳng thức)
giữa các chiều của các lược đồ Picard của $X_{s'}$ và của «chuyên biệt hóa»
$X_s$ của nó. Hơn nữa, ngay từ trước khi các lược đồ Picard được xây dựng,
Serre đã nhận xét rằng tính bất biến của chiều các lược đồ Picard của $X_s$ trong
trường hợp một cấu xạ trơn $f:X\longrightarrow S$ là hệ quả hình thức của lý
thuyết chuyên biệt hóa nhóm cơ bản ([4], X), và của các quan hệ Kummer cổ điển
giữa các điểm có cấp hữu hạn trên đa tạp Picard với nhóm cơ bản Abel hóa ([4],
XI). Nếu nói chung ta ký hiệu bởi $\alpha,\mu,\lambda$ các chiều của phần Abel,
phần nhân và phần cộng của $\operatorname{Pic}_{X_s/k(s)}^\circ$, và định
nghĩa tương tự $\alpha',\mu',\lambda'$, thì các quan hệ đã biết được biểu thị
bởi những bất đẳng thức sau:
\[
\tag{*}\alpha+\mu+\lambda\geq\alpha'+\mu'+\lambda'
\]
(đúng miễn là $\operatorname{Pic}_{X/S}$ tồn tại, do đó có lẽ đúng trong mọi
trường hợp), và khi $\lambda=\lambda'=0$, bất đẳng thức này rút về một đẳng
thức, đúng dưới cùng các giả thiết tồn tại:
\[
\alpha+\mu=\alpha'+\mu' ;
\]
mặt khác ta có
\[
\tag{**}2\alpha+\mu\leq2\alpha'+\mu'
\]
\src{235}
theo lập luận của Serre, nếu các $X_s$ khả ly (thậm chí không cần giả sử
$\operatorname{Pic}_{X/S}^\tau$ tồn tại), hoặc nếu các
${}_n\operatorname{Pic}_{X/S}$ (các hạt nhân của $\varphi_n$ trong hàm tử
Picard) tách trên phần mở của $S$ nơi $n$ khả nghịch (vì chúng étale trên phần mở đó theo (2.5)). Ta có xu hướng
phỏng đoán rằng (*) luôn là một đẳng thức, hoặc ít nhất khi các $X_s$ khả ly,
và rằng ta có các bất đẳng thức
\[
\tag{***}\alpha\leq\alpha',\qquad\lambda\geq\lambda',
\]
những bất đẳng thức này phải đúng mỗi khi ta có một tiền lược đồ nhóm thuộc
kiểu hữu hạn địa phương trên $S$ Noether địa phương, trong đó chiều của các
thớ là hằng (xem (1.3) về một kết quả khẳng định theo hướng này).

\result{NHẬN XÉT}{2.9} Trong mọi trường hợp đã biết,
$\operatorname{Pic}_{X/S}^\tau$ mở phổ dụng trên $S$, nhưng có lẽ không nên vì
thế mà nuôi quá nhiều ảo tưởng, ngay cả khi $f:X\longrightarrow S$ trơn. Dù
sao, Mumford đã xây dựng một ví dụ (đúng là với $S$ không rút gọn, thực ra $S$
là phổ của một vành Artin) trong đó $\operatorname{Pic}_{X/S}^\tau$ không
phẳng trên $S$, bằng cách làm biến thiên vô cùng bé mặt Igusa. Hơn nữa, điểm
được Mumford xét nằm trong $\operatorname{Pic}_{X/S}^\rho$, và vẫn có thể xảy
ra rằng (với $f:X\longrightarrow S$ trơn) $\operatorname{Pic}_{X/S}^\tau$
phẳng trên $S$ tại các điểm của $\operatorname{Pic}_{X/S}^\sigma$; tuy nhiên,
diễn giả nghi ngờ rằng điều này luôn đúng, ngay cả khi giới hạn ở các điểm của
$\operatorname{Pic}_{X/S}^\circ$ và giả sử $S$ là phổ của một vành định giá
rời rạc. Câu hỏi này còn gắn với việc nghiên cứu các điểm bất động của một
lược đồ Abel dưới tác động của một nhóm hữu hạn các tự đẳng cấu, một tình
huống mà dường như ta còn thiếu ví dụ.

Dường như, ngay cả khi chỉ xét $f:X\longrightarrow S$ trơn và xạ ảnh, các kết
quả về tính chính quy địa phương của $\operatorname{Pic}_{X/S}$ được nêu trong
mục này, cùng với các phỏng đoán được chỉ ra trong (2.8), bao quát gần như tất
cả những gì có thể nói về vấn đề này nếu không có các giả thiết cụ thể hơn về bản chất
các thớ của $f$. Tuy nhiên, nhắc lại rằng nếu các thớ hình học của
$\operatorname{Pic}_{X/S}$ rút gọn và không có thành phần cộng, thì từ (1.8)
và (2.5) suy ra rằng $\operatorname{Pic}_{X/S}$ trơn trên $S$ tại các điểm của
$\operatorname{Pic}_{X/S}^\sigma$ khi $S$ rút gọn; kết quả này thậm chí vẫn
đúng, nếu $f:X\longrightarrow S$ chuẩn, mà không có giả thiết nào về $S$, như
ta sẽ thấy trong (3.5). Liên quan đến điều này, ta ghi nhận:

\result{MỆNH ĐỀ}{2.10}

\noindent(i) Nếu $\operatorname{Pic}_{X/S}$ trơn (tương ứng phẳng) trên $S$
tại các điểm của tiết diện đơn vị, thì nó cũng như vậy tại mọi điểm của
$\operatorname{Pic}_{X/S}^\sigma$, và tại các điểm của mọi tiết diện của
$\operatorname{Pic}_{X/S}$ trên $S$.
\src{236}
\noindent(ii) Cho $s\in S$ sao cho
\[
H^2(X_s,\mathcal O_{X_s})=0,
\]
khi đó tồn tại một lân cận mở $U$ của $s$ sao cho
$\operatorname{Pic}_{X/S}|U$ trơn trên $U$.

\noindent(iii) Cho $X$ là một lược đồ riêng trên một trường, khi đó
\[
\dim\operatorname{Pic}_{X/k}\leq\dim H^1(X,\mathcal O_X),
\]
dấu bằng xảy ra khi và chỉ khi $\operatorname{Pic}_{X/k}$ trơn trên $k$; điều
này luôn đúng nếu $k$ có đặc số không.

(i) suy ra từ (2.5) và (1.12), còn (ii) suy ra từ tiêu chuẩn vi phân vô cùng
đối với các cấu xạ trơn và một phép tính chướng ngại quen thuộc, vì (theo
``định lý nửa liên tục'') giả thiết đặt tại $s$ vẫn đúng tại các điểm lân cận.
Cuối cùng, (iii) suy ra từ sự kiện $H^1(X,\mathcal O_X)$ đẳng cấu với không
gian tiếp xúc Zariski của $\operatorname{Pic}_{X/k}$ tại phần tử đơn vị; khẳng
định cuối cùng là một trường hợp riêng của một định lý của Cartier, nói rằng
một ``nhóm hình thức'' trong đặc số 0 là trơn hình thức trên $k$.

\fsection{3}{Lược đồ con Abel chính tắc của
\texorpdfstring{$\operatorname{Pic}_{X/S}$}{Pic(X/S)} và lược đồ Albanese}

\result{MỆNH ĐỀ}{3.1} Cho $k$ là một trường và $G$ là một lược đồ nhóm giao
hoán kiểu hữu hạn trên $k$, ``không có thành phần cộng''. Khi đó
$G_{\mathrm{réd}}^\circ$ khả ly trên $k$, và do đó nó là một lược đồ nhóm con
trơn trên $k$.

Điều này hiển nhiên nếu $k$ hoàn hảo, đặc biệt đối với $G_{\bar k}$, trong đó
$\bar k$ là bao đóng đại số của $k$; vì thế chỉ cần chứng minh rằng
$(G_{\bar k}^\circ)_{\mathrm{réd}}$ được hạ xuống từ một lược đồ con của $G$.
Từ giả thiết trên $G_{\bar k}$ (không có thành phần cộng)\footnote{Bản in ghi
``Bây giờ, giả thiết ...''; đính chính liên quan, tr.~303, quy định ``Từ giả
thiết ...''.} dễ dàng suy ra rằng tồn tại một số nguyên $m$ sao cho
$(G_{\bar k}^\circ)_{\mathrm{réd}}$ là ảnh theo nghĩa lược đồ của đồng cấu lũy
thừa bậc $m$ trong $G_{\bar k}$. Vì đồng cấu cuối này đến từ đồng cấu tương tự
đối với $G$, ảnh theo nghĩa lược đồ của đồng cấu sau đáp ứng yêu cầu.

\result{HỆ QUẢ}{3.2} Cho $X$ là một lược đồ chuẩn hình học và riêng trên $k$ sao cho
$\operatorname{Pic}_{X/k}$ tồn tại. Khi đó tồn tại một lược đồ con Abel $A$
của $\operatorname{Pic}_{X/k}$ có tập nền là
$\operatorname{Pic}_{X/k}^\circ$.

Thật vậy, theo (2.1), (ii), $\operatorname{Pic}_{X/k}^\circ$ riêng trên $k$
nên thỏa mãn điều kiện của (3.1).
\src{237}
Vậy kết quả trước cho thấy rằng, trong một số trường hợp, ``đa tạp Picard''
cổ điển (mà trong lý thuyết hiện nay sẽ được ký hiệu là
$(\operatorname{Pic}_{X/k})_{\mathrm{réd}}^\circ$) ``được xác định trên
$k$'', mà không cần giả sử trường $k$ là hoàn hảo.

Bây giờ, cho $f:X\longrightarrow S$ là một lược đồ tương đối riêng và phẳng,
với
\[
\mathcal O_S\xrightarrow{\sim}f_*(\mathcal O_X)
\]
để đơn giản, sao cho $\operatorname{Pic}_{X/S}$ tồn tại và
$\operatorname{Pic}_{X/S}^\circ$ riêng trên $S$. Ngoài ra, đối với (3.3),
(ii), ta giả sử rằng có một tập mở của $\operatorname{Pic}_{X/S}$ chứa
$\operatorname{Pic}_{X/S}^\circ$ và tựa xạ ảnh trên $S$; như đã thấy, điều
kiện này được thỏa mãn nếu $f$ là xạ ảnh và có các thớ hình học khả ly và bất khả
quy. Nhắc lại rằng một lược đồ Abel trên $S$ là một lược đồ nhóm trên $S$,
riêng và trơn trên $S$, có các thớ liên thông hình học. Ta xét xem có tồn tại
một lược đồ nhóm con $A$ của $\operatorname{Pic}_{X/S}$, vừa là một lược đồ
Abel vừa có tập nền là $\operatorname{Pic}_{X/S}^\circ$, hay không. Ta vừa
thấy rằng một lược đồ như thế luôn tồn tại nếu $S$ là phổ của một trường. Sau
đây là điều có thể khẳng định trong trường hợp tổng quát đang xét:

\result{ĐỊNH LÝ}{3.3} Dưới các điều kiện trên:

\noindent(i) Nếu tồn tại một lược đồ con Abel của
$\operatorname{Pic}_{X/S}$ có tập nền là
$\operatorname{Pic}_{X/S}^\circ$, thì nó là duy nhất. Do đó, phép tạo thành
nó tương thích với các phép đổi cơ sở.

\noindent(ii) Để một lược đồ con Abel như thế tồn tại, chỉ cần nó tồn tại sau
mọi phép đổi cơ sở $S'\longrightarrow S$, trong đó $S'$ là Artin địa phương;
nếu $S$ là phổ của một vành địa phương, thì thậm chí chỉ cần kiểm tra với
$S'=\operatorname{Spec}(A_n)$, trong đó
$A_n=A/\mathfrak m^{n+1}$. Nếu $S$ rút gọn, cũng chỉ cần kiểm tra với các phép
đổi cơ sở $S'\longrightarrow S$, trong đó $S'$ là phổ của một vành định giá
rời rạc (đầy đủ với trường thặng dư đóng đại số, nếu muốn).

\noindent(iii) Giả sử $A$ tồn tại, và cho
$B=\operatorname{Alb}^0(X/S)$ là lược đồ Abel đối ngẫu (tức là
$B=\operatorname{Pic}_{A/S}^\circ$ [7]). Khi đó, ta có thể xây dựng chính tắc
một không gian thuần nhất chính $P=\operatorname{Alb}^1(X/S)$ dưới $B$, và
một $S$-cấu xạ $X\longrightarrow P$ phổ dụng đối với các $S$-cấu xạ từ $X$
vào các lược đồ para-Abel (tức là các không gian thuần nhất chính dưới các
lược đồ Abel). Phép tạo thành $\operatorname{Alb}^0(X/S)$,
$\operatorname{Alb}^1(X/S)$ và
$X\longrightarrow\operatorname{Alb}^1(X/S)$ giao hoán với phép đổi cơ sở.

Ta phác thảo chứng minh:

(i) là một tính chất cứng tổng quát của các lược đồ con Abel trong các lược đồ
nhóm giao hoán: nếu hai lược đồ con như thế trùng nhau theo nghĩa tập hợp tại
một điểm $s\in S$, thì chúng trùng nhau trên toàn bộ thành phần
\src{238}
liên thông của $s$ [7]. (Kết quả này khái quát một định lý cổ điển của Chow.)

(ii) Dùng các lược đồ Hilbert, ta thấy hàm tử gán cho mỗi $S'$ trên $S$ tập
(chỉ gồm 1 hoặc 0 phần tử) các lược đồ con Abel chính tắc của
$(\operatorname{Pic}_{X/S})\times_S S'$ được biểu diễn bởi một lược đồ $T$
kiểu hữu hạn trên $S$. Theo (i), $T\longrightarrow S$ là một đơn cấu, và theo
(3.2), nó là toàn ánh. Nói rằng tồn tại một lược đồ con Abel chính tắc của
$\operatorname{Pic}_{X/S}$ có nghĩa là $T$ có một tiết diện trên $S$, hay
tương đương, $T\longrightarrow S$ là một đẳng cấu. Điều này cũng tương đương
với việc $T\longrightarrow S$ là étale, hoặc, trong trường hợp $S$ rút gọn,
với việc $T\longrightarrow S$ là riêng. Từ đó suy ra ngay (ii).

(iii) Chỉ cần dùng định nghĩa của $\operatorname{Pic}_{X/S}$, ta thấy rằng,
đối với mọi lược đồ Abel $C$ trên $S$, việc cho một $S$-cấu xạ từ $X$ vào một
không gian thuần nhất chính dưới $C$ tương đương với việc cho một đồng cấu nhóm
$C'\longrightarrow\operatorname{Pic}_{X/S}$, trong đó $C'$ là lược đồ Abel
đối ngẫu của $C$. Mà nếu lược đồ con Abel chính tắc $A$ của
$\operatorname{Pic}_{X/S}$ tồn tại, thì các đồng cấu này nhất thiết phân tích
qua $A$ (cũng có thể thấy điều này bằng cách dùng các điểm có cấp hữu hạn). Từ
đó suy ra ngay (iii).

\result{NHẬN XÉT}{3.4} Ta sẽ ký hiệu bởi
$\operatorname{Pic}_{X/S}^{\circ\circ}$ lược đồ con Abel chính tắc của
$\operatorname{Pic}_{X/S}$, nếu nó tồn tại. Đáng tiếc là điều này không phải
lúc nào cũng xảy ra, như có thể thấy khi làm biến thiên vô cùng bé mặt Igusa
(bằng biến dạng modular bậc nhất). Tuy nhiên, có thể
$\operatorname{Pic}_{X/S}^{\circ\circ}$ ít nhất tồn tại nếu $S$ rút gọn,
hoặc, theo (ii), tương đương với điều đó, nếu $S$ là phổ của một vành định giá
rời rạc. Khi đó, cho $X_0,X_1$ là các thớ\footnote{Bản in dùng danh từ số ít
``la fibre'' cho hai thớ; cách đọc số nhiều là ``les
fibres''.} đặc biệt và tổng quát của $X$, và đặt
\[
A_1=\operatorname{Pic}_{X_1/K}^\circ,
\]
trong đó $K$ là trường phân thức của vành định giá $V$. Theo Koizumi, tồn tại
một lược đồ Abel $A$ trên $S$, duy nhất về căn bản, có thớ tổng quát là $A_1$,
và dễ dàng thấy như trong (2.1), (i) (từ đây giả sử $X$ trơn trên $S$), rằng
cấu xạ đồng nhất của $A_1$ mở rộng thành một cấu xạ
\[
A\longrightarrow\operatorname{Pic}_{X/S}.
\]
Do đó có một đồng cấu
\[
\tag{*}A_0\longrightarrow\operatorname{Pic}_{X_0/k}^{\circ\circ},
\]
mà ta dễ dàng chứng minh là toàn ánh, có hạt nhân là một $p$-nhóm hữu hạn,
trong đó $p$ là số mũ đặc số của trường thặng dư $k$
\src{239}
(vẫn bằng cách dùng các điểm có cấp hữu hạn). Với quy ước đó, các điều kiện
sau đây trên $X/S$ là tương đương:

\medskip
\noindent a. Đồng cấu (*) ở trên là một đẳng cấu (điều mà Shimura sẽ diễn đạt
bằng cách nói rằng phép tạo thành ``đa tạp Picard'' là ``tương thích với phép
chuyên biệt hóa'').

\noindent b. $\operatorname{Pic}_{X/S}^{\circ\circ}$ tồn tại (khi đó nó
không gì khác ngoài $A$).

\noindent c. (Để ghi nhớ) Các
$\operatorname{Pic}_{X_n/S}^{\circ\circ}$ tồn tại.

Theo nhận xét đã nêu về hạt nhân của (*), điều kiện (a) được thỏa mãn nếu đặc
số thặng dư bằng không, nhưng kết quả này sẽ được khái quát đáng kể trong
(3.5).

Dĩ nhiên, nếu $\operatorname{Pic}_{X/S}$ trơn trên $S$ tại các điểm của
$\operatorname{Pic}_{X/S}^\circ$, thì tập sau là mở trong
$\operatorname{Pic}_{X/S}$ (xem (1.7)) và do đó, với cấu trúc cảm sinh, nó là
một lược đồ con Abel của $\operatorname{Pic}_{X/S}$, nên bằng
$\operatorname{Pic}_{X/S}^{\circ\circ}$, vốn tồn tại trong trường hợp này.
Nhưng ta còn có kết quả mạnh hơn nhiều:

\result{ĐỊNH LÝ}{3.5} Dưới các điều kiện của (3.3), cho $s\in S$ sao cho
$\operatorname{Pic}_{X_s/k(s)}$ trơn trên $k(s)$ (hay, tương đương,
\[
\dim\operatorname{Pic}_{X_s/k(s)}=\dim H^1(X_s,\mathcal O_{X_s})).
\]
Khi đó tồn tại một lân cận mở $U$ của $s$ sao cho
$\operatorname{Pic}_{X/S}$ trơn trên $S$ tại các điểm của\footnote{Bản in ghi
cấu tạo bất chuẩn ``en les points de''; cách viết chuẩn là ``aux
points de''.} $\operatorname{Pic}_{X/S}^\circ|U$, và do đó đây là một lược
đồ con Abel mở trong $\operatorname{Pic}_{X/S}|U$. Đặc biệt,
$\operatorname{Pic}_{X|U/U}^{\circ\circ}$ tồn tại.

Ta nêu nguyên lý của chứng minh. Những điều trên đưa ta về trường hợp $S$ là
phổ của một vành artin địa phương $A$, và ta quy nạp theo cấp vô cùng bé của
$A$. Do đó, có thể giả sử $\operatorname{Pic}_{X_n/A_n}^\circ$ trơn trên
$A_n$, và chỉ cần chứng minh rằng
$\operatorname{Pic}_{X_{n+1}/A_{n+1}}^\circ$ trơn trên $A_{n+1}$. Ta thấy
rằng để làm điều này, chỉ cần xây dựng một lược đồ Abel $P_{n+1}$ trên
$A_{n+1}$ mở rộng
\[
P_n=\operatorname{Pic}_{X_n/A_n}^\circ,
\]
đồng thời xây dựng một Môđun khả nghịch $\mathcal L_{n+1}$ trên
$X_{n+1}\times_{A_{n+1}}P_{n+1}$ mở rộng Môđun khả nghịch $\mathcal L_n$
trên $X_n\times_{A_n}P_n$ xuất hiện trong định nghĩa của lược đồ Picard
$\operatorname{Pic}_{X_n/A_n}$ như nghiệm của một bài toán phổ dụng. (N.B. ---
Ta có thể giả sử $X$ có một tiết diện trên $S$~...) Để thực hiện phép xây dựng,
phải dùng kết quả then chốt sau: mọi lược đồ Abel xác định trên một vành thương
của một vành artin địa phương đều có thể mở rộng (nói cách khác, ``lược đồ
hình thức các Môđun'' tuyệt đối ([3], II) của một lược đồ Abel trên một trường
đóng đại số là trơn trên vành các
\src{240}
vectơ Witt của $k$); bản thân kết quả này thu được bằng cách dùng các tính chất
hình thức tổng quát của chướng ngại đối với phép nâng và các phép toán nhóm.
Chấp nhận kết quả này, trước hết ta mở rộng $P_n$ một cách tùy ý thành
$P_{n+1}$; khi đó gặp một chướng ngại đối với việc nâng $\mathcal L_n$, nằm
trong
\[
H^2(X_0\times P_0,\mathcal O_{X_0\times P_0})\otimes_k V
\qquad(\text{trong đó }V=\mathfrak m^{n+1}/\mathfrak m^{n+2}),
\]
và thực ra nằm trong không gian con
\[
H^1(X_0,\mathcal O_{X_0})\otimes
H^1(P_0,\mathcal O_{P_0})\otimes_k V
\]
(do hạn chế của $\mathcal L_n$ lên hai nhân tử $X_n$ và $P_n$ là tầm thường).
Mà không gian sau chính là
\[
H^1(P_0,\mathcal T_{P_0/k})\otimes V,
\]
trong đó $\mathcal T_{P_0/k}$ là bó tiếp xúc của $P_0/k$, vì vậy nó cũng là
không gian biểu thị tính bất định trong phép nâng $P_n$ thành $P_{n+1}$ ([3],
II). Do đó, có thể hiệu chỉnh phép nâng này (theo một và chỉ một cách, như phải
vậy) để triệt tiêu chướng ngại đối với việc nâng $\mathcal L_n$.

\result{HỆ QUẢ}{3.6} Dưới các điều kiện của (3.5),
$R^1f_*(\mathcal O_X)$ là một Môđun tự do địa phương trên $S$ trong một lân
cận của $s$, và phép tạo thành nó giao hoán với phép đổi cơ sở.

Thật vậy, Môđun này chính là Môđun tiếp xúc của
$\operatorname{Pic}_{X/S}$ dọc theo tiết diện đơn vị.

\result{NHẬN XÉT}{3.7} Dùng cùng lập luận như đối với (3.5), ta có thể chứng
minh rằng nếu $S'$ là một lược đồ con của $S$ được xác định bởi một Iđêan
coherent lũy linh, và nếu chỉ giả sử $\operatorname{Pic}_{X'/S'}^\circ$ tồn
tại và có các thớ trơn, thì $\operatorname{Pic}_{X/S}^\circ$ nhất thiết tồn
tại và là một lược đồ Abel trên $S$. Điều này cho phép xây dựng các lược đồ
Picard trong một số trường hợp mặc dù không có giả thiết xạ ảnh; chẳng hạn,
lược đồ Abel đối ngẫu của một lược đồ Abel trên một vành artin luôn tồn tại.
Mặt khác, dùng (3.6) trong trường hợp $X$ là một lược đồ Abel trên $S$, và dùng
cấu trúc đã biết của $H^*(X_s,\mathcal O_{X_s})$ như một đại số ngoài trên
$H^1(X_s,\mathcal O_{X_s})$ (Rosenlicht--Serre), ta thấy rằng với mọi $i$,
$R^if_*(\mathcal O_X)$ tự do địa phương, và chính xác hơn, đẳng cấu với lũy
thừa ngoài thứ $i$ của $R^1f_*(\mathcal O_X)$.

\result{NHẬN XÉT}{3.8} Trong trường hợp một cấu xạ xạ ảnh trơn
$f:X\longrightarrow S$, với $S$ rút gọn và có các đặc số thặng dư bằng không,
kết quả (3.6) đã được biết bằng phương pháp siêu việt, như một hệ quả của lý
thuyết Hodge. Thật ra, tất cả các $R^pf_*(\Omega^q_{X/S})$ khi đó đều tự do
địa phương. Trái lại, có các phản ví dụ
\src{241}
trong trường hợp đặc số hỗn hợp đối với tính tự do địa phương của
$R^1f_*(\mathcal O_X)$, bằng cách dùng các đa tạp Serre ([8], tr.~50). Dường
như không có phản ví dụ nào trong trường hợp cùng đặc số.

\result{HỆ QUẢ}{3.8} Dưới các điều kiện của (3.5),
$\operatorname{Pic}_{X/S}|U$ trơn trên $U$ tại mọi điểm của
$\operatorname{Pic}_{X/S}^\sigma|U$.

Áp dụng (2.10), (i). Đặc biệt, có tính đến (2.10), (iii):

\result{HỆ QUẢ}{3.9} Giả sử $S$ có các đặc số thặng dư bằng không. Khi đó
$\operatorname{Pic}_{X/S}^\tau$ trơn trên $S$.

\result{NHẬN XÉT}{3.10} Từ đó suy ra, chẳng hạn, nếu
$\operatorname{Pic}_{X/S}^\tau$ cũng riêng trên $S$, thì nhóm con xoắn của các
nhóm Néron--Severi của các thớ hình học của $f$ không đổi trên mỗi thành phần
liên thông của $S$ (điều này cũng hiển nhiên bằng phương pháp siêu việt khi
$f$ trơn và xạ ảnh). Ta lưu ý rằng việc áp dụng trực tiếp (2.5) cho phép chứng
minh tổng quát hơn rằng, nếu $\operatorname{Pic}_{X/S}^\tau$ riêng trên $S$ (chẳng
hạn $f:X\longrightarrow S$ trơn và xạ ảnh) và nếu $q$ là một số nguyên tố
khác các đặc số thặng dư của $S$, thì thành phần $q$-nguyên sơ của các nhóm
con xoắn Néron--Severi của các thớ hình học của $X/S$ không đổi trên mỗi thành
phần liên thông của $S$. Tuy nhiên, trong đặc số $p>0$, điều này không còn đúng
đối với thành phần $p$-nguyên sơ của phần xoắn. Dù vậy, vẫn có thể xảy ra rằng
hạng trên $k(s)$ của
\[
\operatorname{Pic}_{X_s/k(s)}^\tau/
\operatorname{Pic}_{X_s/k(s)}^{\circ\circ}=T_{X_s/k(s)}
\]
là hằng địa phương; khi $S$ rút gọn, có thể chứng minh rằng điều này tương
đương với việc $\operatorname{Pic}_{X/S}^{\circ\circ}$ tồn tại và
$\operatorname{Pic}_{X/S}^\tau$ phẳng trên $S$, và chỉ cần kiểm tra câu hỏi
khi $S$ là phổ của một vành định giá rời rạc. Đó là điều tôi đã kiểm tra trong
một vài ví dụ đã xét; nhưng vì phát biểu tương ứng với $S$ artin là sai (xem
các nhận xét (2.9) và (3.4)), có lẽ không nên nuôi quá nhiều ảo tưởng.

\fsection{4}{Định lý về tính hữu hạn đối với lược đồ Picard}

Cho $f:X\longrightarrow S$ là một cấu xạ xạ ảnh và phẳng sao cho
$\operatorname{Pic}_{X/S}$ tồn tại; khi đó, việc xét các ``đa thức Hilbert''
$Q\in\mathbf Q[t]$ cho phép phân tích $\operatorname{Pic}_{X/S}$ thành tổng
rời nhau của các tập mở $\operatorname{Pic}_{X/S}^Q$. Nếu không đặt thêm các
giả thiết, chẳng hạn bảo đảm $\operatorname{Pic}_{X/S}$ tách trên $S$, thì
nói chung các tập mở này không nhất thiết thuộc kiểu hữu hạn trên $S$; ta có
một phản ví dụ khi $X$ là ``một đường conic suy biến thành hai đường thẳng''.
Điều này
\src{242}
có thể vẫn đúng nếu $f$ khả ly và có các thớ hình học bất khả quy. Câu hỏi này
liên quan đến việc $\operatorname{Pic}_{X/S}^\tau$ có thuộc kiểu hữu hạn trên
$S$ hay không, điều có lẽ đúng mà không cần giả thiết nào về các thớ của
$X/S$. Khi $f:X\longrightarrow S$ trơn, lưu ý rằng
$\operatorname{Pic}_{X/S}^\tau$ nằm trong một trong các
$\operatorname{Pic}_{X/S}^Q$ (nhờ sự kiện rằng, trên một đa tạp xạ ảnh không
kỳ dị, tương đương ``theo xoắn'' tinh hơn tương đương số đối với các ước, và
thật ra, theo Matsusaka, là ``đồng nhất với'' tương đương số), nên nó thuộc
kiểu hữu hạn trên $S$ nếu các $\operatorname{Pic}_{X/S}^Q$ cũng vậy. Lưu ý
rằng các câu hỏi về tính hữu hạn vừa nêu vẫn có một ý nghĩa hiển nhiên ngay cả khi
không giả sử $\operatorname{Pic}_{X/S}$ tồn tại, vì chúng được diễn đạt bằng
cách nói rằng một số họ Môđun khả nghịch là ``bị chặn'' theo nghĩa của [3],
IV: với mỗi trường đóng đại số $k$, xét các lược đồ con nguyên của
$\mathbf P_k^r$, có chiều $n$ và bậc $d$, cùng các Môđun khả nghịch trên các
tiền lược đồ như thế, có đa thức Hilbert $Q$ (trong đó $r,n,d,Q$ được cho
trước); cần chứng minh rằng họ các Môđun này (được xem là các Môđun coherent
trên các $\mathbf P_k^r$) bị chặn, tức là có thể được tham số hóa bởi một lược
đồ kiểu hữu hạn trên $\mathbf Z$.

Dùng phương pháp của Matsusaka [6], một chứng minh khá kỹ thuật (viện đến các
``tiêu chuẩn tương đương'' dưới một dạng thích hợp) cho phép trả lời khẳng định
khi ta giới hạn ở các đa tạp con không kỳ dị của $\mathbf P_k^r$. Chính xác
hơn, ta được kết quả sau:

\result{ĐỊNH LÝ}{4.1} Cho $f:X\longrightarrow S$\footnote{Bản in ghi
``Soient : $X\longrightarrow S$'' (dạng số nhiều); cấu xạ được dùng về sau là
$f$.} là một cấu xạ trơn và xạ ảnh có các thớ liên thông hình học, với $S$
noether, $\mathcal O_X(1)$ là một Môđun rất rộng trên $X$, tương đối với $S$,
$E$ là một tập con của $\operatorname{Pic}_{X/S}$ tương ứng với một họ
$(\mathcal L_i)$ các Môđun khả nghịch trên các thớ hình học
$\overline X_{s_i}$ của $X/S$, $D_i$ là một ước (không nhất thiết dương) trên
$\overline X_{s_i}$ xác định $\mathcal L_i$,
\[
a_n^{(i)}t^n+\cdots+a_0^{(i)}
\]
là đa thức Hilbert của $\mathcal L_i$, $\xi=\xi_i$ là một ước xác định
$\mathcal O_X(1)$, tức là một tiết diện siêu phẳng. Các điều kiện sau là tương
đương:

\noindent a. $E$\footnote{Bản in ghi tại đây ``$Q$''; tập con được đưa vào
trong phát biểu và có tính tựa compact biểu thị họ bị chặn là $E$.} là tựa
compact, tức là nằm trong một tập mở kiểu hữu hạn trên $S$, tức là họ các
$(\mathcal L_i)$ bị chặn.

\noindent b. $E$ nằm trong hợp của một số hữu hạn các tập
$\operatorname{Pic}_{X/S}^Q$, $Q\in\mathbf Q[t]$, tức là tập các đa thức
Hilbert của các $\mathcal L_i$ là hữu hạn.

\noindent b'. (Nếu các thớ của $X/S$ đều có cùng chiều $n$.) Hai hệ số
$a_{n-1}^{(i)}$ và $a_{n-2}^{(i)}$ của các đa thức Hilbert của
$\mathcal L_i$ nằm trong một tập hữu hạn.
\src{243}
\result{HỆ QUẢ}{4.2} Cho $f:X\longrightarrow S$ là một cấu xạ trơn và xạ ảnh
có các thớ liên thông hình học; khi đó các lược đồ
$\operatorname{Pic}_{X/S}^Q$ và $\operatorname{Pic}_{X/S}^\tau$
($Q\in\mathbf Q[t]$) là xạ ảnh trên $S$.

Vì chúng thuộc kiểu hữu hạn trên $S$ theo (4.1), ta có thể áp dụng (2.4).

\section*{TÀI LIỆU THAM KHẢO}
\addcontentsline{toc}{section}{Tài liệu tham khảo}

\begin{description}
\item[{[1]}] CHEVALLEY (Claude). -- Sur la théorie de Picard,
\emph{Amer. J. of Math.}, t. 82, 1960, tr. 435--490.

\item[{[2]}] GROTHENDIECK (A.) và DIEUDONNÉ (J.). -- Éléments de géométrie
algébrique, Chương I và các chương tiếp theo. -- Paris, Presses universitaires
de France, 1960--1961 (Institut des Hautes Études Scientifiques, 4, 8, 11,
\ldots).

\item[{[3]}] GROTHENDIECK (Alexander). -- Technique de descente et théorèmes
d'existence en géométrie algébrique, I--V, \emph{Séminaire Bourbaki}, t. 12,
1959/60, n\textsuperscript{o} 190, 29 tr. và n\textsuperscript{o} 195, 22 tr.;
t. 13, 1960/61, n\textsuperscript{o} 212, 20 tr. và
n\textsuperscript{o} 221, 28 tr.; t. 14, 1961/62,
n\textsuperscript{o} 232, 19 tr.

\item[{[4]}] GROTHENDIECK (Alexander). -- Séminaire de géométrie algébrique,
1960/61. -- Paris, Institut des Hautes Études Scientifiques, 1961.

\item[{[5]}] GROTHENDIECK (Alexander). -- Séminaire de géométrie algébrique,
1961/62, do một nhóm thính giả biên soạn. -- Paris, Institut des Hautes Études
Scientifiques (sắp xuất bản).

\item[{[6]}] MATSUSAKA (T.). -- The criteria for algebraic equivalence and the
torsion group, \emph{Amer. J. Math.}, t. 79, 1957, tr. 53--66.

\item[{[7]}] MUMFORD (D.) và TATE (J.). -- Séminaire de géométrie algébrique,
Harvard University, học kỳ mùa xuân 1962 (sắp xuất bản).

\item[{[8]}] SERRE (Jean-Pierre). -- Sur la topologie des variétés algébriques
en caractéristique $p$, \emph{Symposium internacional de topología algebraica}
[1956, Mexico], tr. 24--53. -- Mexico, Universidad nacional autónoma, 1958.
\end{description}

\par
